\section*{第 \ref{ch:b3:supramolecular} 章——超分子化学}

\begin{solution}{exo:b3:supramolecular:1}
离子–偶极；$\pi$ 堆积；阳离子--$\pi$；卤键（I 到 N）；三个氢键。
\end{solution}

\begin{solution}{exo:b3:supramolecular:2}
$\Delta G^\circ = -RT\ln K = -8.314 \times 298 \times \ln(\num{1.0e4}) = \qty{-22.8}{kJ/mol}$.
\end{solution}

\begin{solution}{exo:b3:supramolecular:3}
$K_{\ce{K+}}/K_{\ce{Na+}} = 10^{1.7} = 50$；自由能差为 $RT\ln 50 = \qty{9.7}{kJ/mol}$，对钾有利。
\end{solution}

\begin{solution}{exo:b3:supramolecular:4}
$n/(n + m) = 1/3$：一个主体对应两个客体，即 $\mathrm{HG_2}$。
\end{solution}

\begin{solution}{exo:b3:supramolecular:5}
$H_0 + G_0 + 1/K = \qty{4.667e-3}{mol/L}$；$[\mathrm{HG}] = (4.667 - \sqrt{4.667^2 - 4 \times 2.0 \times 2.0})/2 \times 10^{-3} = \qty{1.13e-3}{mol/L}$：结合
分数 0.566，$\delta = 3.560 + 0.160 \times 0.566 = \qty{3.651}{ppm}$。
\end{solution}

\begin{solution}{exo:b3:supramolecular:6}
$\log K = 18.3 - 8.6 = 9.7$，$K = \num{5e9}$；$\Delta G^\circ = -RT\ln K = \qty{-55}{kJ/mol}$。六个 Ni--N 键
两侧属同一类，游离分子数由四个增至
七个：平动熵的增加驱动反应进行。
\end{solution}

\begin{solution}{exo:b3:supramolecular:7}
$\Delta G^\circ = -RT\ln(\num{2.0e5}) = \qty{-30.2}{kJ/mol}$; $T\Delta S^\circ = \Delta H^\circ - \Delta G^\circ = -40 + 30.2 = \qty{-9.8}{kJ/mol}$;
$\Delta S^\circ = \qty{-33}{J.K^{-1}.mol^{-1}}$：焓驱动的结合，部分由熵偿还。
\end{solution}

\begin{solution}{exo:b3:supramolecular:8}
铜(I) 与两个双齿菲啰啉单元结合，并因呈四面体构型，
使二者成直角排列，其中一条穿过另一条将闭合成环的空间，
两链闭合成环后即互锁。加入大量过量的氰根离子，
它与铜(I) 的结合强得多，可夺去金属，
留下游离的索烃。
\end{solution}

\begin{solution}{exo:b3:supramolecular:9}
$k_c = \pi\Delta\nu/\sqrt2 = \pi \times 120/1.414 = \qty{267}{s^{-1}}$。由艾林方程：
$\Delta G^\ddagger = RT\ln(k_BT/hk_c) = 8.314 \times 300 \times \ln(\num{6.25e12}/267) = \qty{59.6}{kJ/mol}$.
\end{solution}

\begin{solution}{exo:b3:supramolecular:10}
药物过量且为固体时，游离药物的浓度固定为 $S_0$。于是 $[\mathrm{DC}] = KS_0[\mathrm C]$，并且
$C_t = [\mathrm C] + [\mathrm{DC}] = [\mathrm C](1 + KS_0)$，故 $[\mathrm{DC}] = KS_0C_t/(1 + KS_0)$，总溶解度为 $S_0 + [\mathrm{DC}]$。这里
$KS_0 = 0.050$：$S = 0.10 + 0.050 \times 10/1.050 = \qty{0.58}{mmol/L}$，接近内在值的六倍。
\end{solution}

\begin{solution}{exo:b3:supramolecular:11}
当 $1/K \ll H_0, G_0$ 时，$[\mathrm{HG}] \approx \bigl(H_0 + G_0 - |H_0 - G_0|\bigr)/2 = \min(H_0, G_0)$：结合分数为 $G_0/H_0$，直至一
当量，之后为 1。曲线不再依赖于 $K$：任何足够大的值
在误差范围内给出相同的尖锐转折，因此数据只能说明 $K$ 超过
某一下限。
\end{solution}

\begin{solution}{exo:b3:supramolecular:12}
$k_{\mathrm{off}} = k_{\mathrm{on}}/K = \qty{1e9}{L.mol^{-1}.s^{-1}}/\qty{1500}{L.mol^{-1}} = \qty{6.7e5}{s^{-1}}$。化学位移差
$\qty{0.160}{ppm}$ 在 \qty{400}{MHz} 下为 \qty{64}{Hz}；合并需要 $k \approx \pi \times 64/\sqrt2 = \qty{140}{s^{-1}}$。交换
快数千倍，只出现一个平均信号。
\end{solution}

\begin{solution}{pb:b3:supramolecular:1}
\textbf{1.} 由六个 $\ce{-CH2CH2O-}$ 单元组成的环；配合物中 $\ce{K+}$ 位于中心，由
朝内的六个氧原子结合，$\ce{CH2}$ 基团位于外侧。

\textbf{2.} $\ce{K+}$: $-RT\ln 10 \times 6.1 = \qty{-34.8}{kJ/mol}$; $\ce{Na+}$: \qty{-25.1}{kJ/mol}.

\textbf{3.} $10^{6.1 - 4.4} = 50$.

\textbf{4.} 较大的 $\ce{K+}$（\qty{138}{pm}）能同时接触环上的全部六个氧原子；
较小的 $\ce{Na+}$（\qty{102}{pm}）则不能，除非环发生折叠，但折叠会带来张力。

\textbf{5.} 其外侧为烃类表面，电荷被包埋；$\ce{KCl}$
需要其离子被溶剂化，而甲苯无法做到。

\textbf{6.} 水对 $\ce{K+}$ 的溶剂化远强于甲醇，而离子进入环时必须失去
这部分水才能进入环。

\textbf{7.} $\ce{K+}(\text{水相}) + \ce{MnO4-}(\text{水相}) + \mathrm L(\text{有机相}) \rightleftharpoons \ce{[KL]+MnO4-}(\text{有机相})$。

\textbf{8.} $[\ce{K+}]_{\mathrm{aq}}$ 保持为 \qty{0.10}{mol/L}：$y = \num{1.0e5} \times 0.10 \times (\num{1.0e-3} - y)^2 = \num{1.0e4}(\num{1.0e-3} - y)^2$。

\textbf{9.} 令 $u = \num{1.0e-3} - y$，则 $\num{1.0e4}u^2 + u - \num{1.0e-3} = 0$，$u = \qty{2.70e-4}{mol/L}$，$y = \qty{7.30e-4}{mol/L}$：73\,\%
被萃取。

\textbf{10.} $y = \num{1.0e4}(\num{1.0e-3} - y)(\num{1.0e-2} - y)$ 给出 $y = \qty{9.89e-4}{mol/L}$：99\,\%。

\textbf{11.} $y = \num{1.0e4}(\num{1.0e-3} - y)(\num{1.0e-4} - y)$ 给出 $y = \qty{9.0e-5}{mol/L}$：9\,\%，受限于冠醚，
其负载率为 90\,\%。

\textbf{12.} 高锰酸根离子在甲苯中保持紫色。在其中，它
未被溶剂化，是裸露的阴离子，其反应速率远快于水中；
水中则由水合壳层对其加以屏蔽。

\textbf{13.} 钠离子的结合与解离远快于两共振频率之
差：快速交换。

\textbf{14.} $\delta = \delta_{\mathrm{free}} + (\delta_{\mathrm{bound}} - \delta_{\mathrm{free}})[\mathrm{HG}]/H_0$，其中 $[\mathrm{HG}]$ 取自精确等温线，$H_0 = \qty{2.0e-3}{mol/L}$，
$G_0 = (\text{当量})\,H_0$，$\delta_{\mathrm{free}} = \qty{3.560}{ppm}$。

\textbf{15.} $K = (1.55 \pm 0.08) \times 10^3\ \unit{L.mol^{-1}}$, $\delta_{\mathrm{bound}} = \qty{3.719 \pm 0.001}{ppm}$.

\textbf{16.} 残差的方均根为 \qty{0.0014}{ppm}，与读数精密度相当，
且无明显趋势：1:1 模型拟合良好。

\textbf{17.} 结合分数 $f$ 由 $G_0 = fH_0 + f/(K(1 - f))$ 给出：$f = 0.2$ 时为 0.28 当量，$f = 0.8$ 时为 2.1 当量。

\textbf{18.} $KH_0 = 10^{6.1} \times \num{2.0e-3} \approx 2500$：曲线将在一当量处尖锐转折，
只能给出下限，需要在较低浓度下进行
竞争实验或量热法。

\textbf{19.} 在甲苯中，烯烃与高锰酸根形成均匀溶液。

\textbf{20.} 被还原的锰以二氧化锰形式离开有机相，
而冠醚在界面处释放，再结合另一个 $\ce{K+}$ 和 $\ce{MnO4-}$。每个冠醚
分子可多次携带高锰酸根，只需催化量即可。

\textbf{21.} $K_{\mathrm{ex}}$ 乘以 $[\ce{K+}]$：较高的钾浓度推动萃取进行。

\textbf{22.} 季铵盐，例如四丁基铵盐或
氯化苄基三乙基铵，其阳离子本身即具亲脂性。

\textbf{23.} 苯可致白血病；甲苯能起同样作用，而危害
小得多。

\textbf{24.} 冠醚为 \qty{1.0}{mmol/L} 时，高锰酸根的 73\,\% 被萃取
到甲苯中。
\end{solution}
